\documentclass[a4paper,14pt,oneside,final]{report}
\input{preamble.tex}

\begin{document}

\begin{titlepage}
  \begin{center}
    Министерство образования Республики Беларусь\\[1em]
    Учреждение образования\\
    БЕЛОРУССКИЙ ГОСУДАРСТВЕННЫЙ УНИВЕРСИТЕТ \\
    ИНФОРМАТИКИ И РАДИОЭЛЕКТРОНИКИ\\[1em]

    \begin{minipage}{\textwidth}
      \begin{flushleft}
        \begin{tabular}{ l l }
          Факультет & информационных технологий и управления\\
          Кафедра   & интеллектуальных информационных технологий
        \end{tabular}
      \end{flushleft}
    \end{minipage}\\[1em]

    \begin{flushright}
      \begin{minipage}{0.4\textwidth}
        \textit{К защите допустить:}\\[0.8em]
        Заведующий кафедрой ИИТ\\[0.45em]
        \underline{\hspace*{2.8cm}} Д.\,В.~Шункевич
      \end{minipage}\\[2.2em]
    \end{flushright}

    {ПОЯСНИТЕЛЬНАЯ ЗАПИСКА}\\
    {к курсовой работе
    по дисциплине \\
    «Технологии и инструментальные средства проектирования интеллектуальных систем»}\\
    {на тему:}\\[1em]
    \textbf{\large ИНТЕЛЛЕКТУАЛЬНЫЙ МЕДИЦИНСКИЙ АССИСТЕНТ}\\[1em]


    {БГУИР КР7 1-40 03 01 073 ПЗ}\\[2em]
    %1-40 03 01 01 если специализация ИГС
    %1-40 03 01 02 если специализация ИКТЗИ
    % последние 3 цифры зачетки
    \vspace{2em}
    \begin{tabular}{ p{0.65\textwidth}p{0.25\textwidth} }
      Студент гр. 221701 & И.\,Д.~Телица \\
      Руководитель & Д.\,А.~Сальников \\
    \end{tabular}
    
    \vfill
    {\normalsize Минск 2025}
  \end{center}
\end{titlepage}


%\input{cource/course_task.tex}

%\input{practice/sec_refer.tex}

%\pagenumbering{arabic}
%\setcounter{page}{5}

\input{table_of_contents}

\input{diploma/diploma_content}

%\input{references}


[1] Woebot : [сайт]. — URL: \url{https://woebothealth.com} (дата обращения: 25.11.2025)

[2] HealthTap : [сайт]. — URL: \url{https://www.healthtap.com} (дата обращения: 25.11.2025)

[3] Ada Health : [сайт]. — URL: \url{https://ada.com} (дата обращения: 25.11.2025)

[4] Semantic Web : [сайт]. — URL: \url{https://www.ontotext.com/knowledgehub/fundamentals/what-is-the-semantic-web} (дата обращения: 25.11.2025)

[5] Голенков В. В., Гулякина Н. А., Шукевич Д. В.  
Открытая технология онтологического проектирования, производства и эксплуатации семантически совместимых гибридных интеллектуальных компьютерных систем. — Минск : БГУИР, 2021. — 312 с.

[6] Ralston S., Penman I. D., Strachan M. W. J.  
Davidson’s Principles and Practice of Medicine. — New York : McGraw Hill, 2022. — 4384 p.

[7] Шукевич Д. В., Голенков В. В.  
Агентно-ориентированные модели, методика и средства разработки совместных решателей задач интеллектуальных систем [сайт] // Программные продукты и системы. — 2020. — URL: \url{https://libeldoc.bsuir.by/bitstream/123456789/42228/1/Golenkov_Agentno_orientirovannyye.pdf} (дата обращения: 25.11.2025)

[8] Гулякина Н. А., Давыденко И. Т.  
Семантические модели и метод согласованной разработки баз знаний [сайт] // Программные продукты и системы. — 2020. — URL: \url{https://cyberleninka.ru/article/n/semanticheskie-modeli-i-metod-soglasovannoy-razrabotki-baz-znaniy} (дата обращения: 25.11.2025)

[9] Профилактическая медицина : [сайт]. — URL: \url{https://www.mediasphera.ru/issues/profilakticheskaya-meditsina/2022/2/downloads/ru/1230549482022021081} (дата обращения: 25.11.2025)

[10] Общее устройство человека и его становление : [сайт]. — URL: \url{https://cyberleninka.ru/article/n/obschee-ustroystvo-cheloveka-i-ego-stanovlenie} (дата обращения: 25.11.2025)



% Если в тексте записки есть приложения, для добавления их в содержание, необходимо оставить следующую строку
% Если приложений нет - данную строку добавляем в комментарий
%\include{sections/appendix}


\end{document}